% ---------------------------------------------------------------------------
%  06-theory.tex -- the statements and proofs that are genuinely ours
% ---------------------------------------------------------------------------
\section{Theory: what has to be proved, and the proofs}
\label{sec:theory}

This section states and proves the properties the design relies on.  Results
imported from the literature were stated in \cref{sec:background}
(\cref{prop:alu-finite,prop:extra-sound}) and are used, not reproved.
Everything below is proved here.

Throughout, fix a discrete state $\dk$ with its bounds $(L,U)$, and write
$A$ for the antichain stored at $\dk$.  Recall (\cref{lem:preorder}) that
$\subs$ is a preorder for a fixed $(L,U)$, and that $(L,U)$ is constant within
an antichain.  Define the \emph{covered region}
\[
  \Cov(A) \;=\;
  \begin{cases}
    \bigcup_{Z \in A} \aLU([Z]) & \text{in subsumption mode,}\\[2pt]
    \bigcup_{Z \in A} [Z]       & \text{in exact mode.}
  \end{cases}
\]
Both cases are of the form $\bigcup_{Z\in A} \kappa([Z])$ for a closure
operator $\kappa$ ($\aLU$, or the identity), which is all the proofs use.

\subsection{The concurrency argument}

\begin{lemma}[Replacement]\label{lem:replacement}
Every execution of \code{add} that retires a stored zone $Z$ also inserts, at
the same discrete state and stamped with the current BFS level, a zone $z^\star$
with $Z \subs z^\star$.
\end{lemma}

\begin{proof}
Retirement happens only in steps (3) and (4) of \cref{sec:antichain}, and both
are on the path that reaches step (5): step (2) is the last point at which
\code{add} can return early, and once it is passed, step (5) executes
unconditionally.

Let $z_0$ be the incoming zone and $z^\star$ the zone actually stored in step
(5).  Step (4) only ever \emph{grows} the pending zone --- it replaces $z$ by the
convex hull $z \sqcup Z_i$, and $[z] \subseteq [z \sqcup Z_i]$ --- so
$[z_0] \subseteq [z^\star]$ and, for every $Z_i$ absorbed in step (4),
$[Z_i] \subseteq [z^\star]$.  By \cref{lem:incl-implies-alu}, in both cases
$Z_i \subs z^\star$.

For a zone $Z$ retired in step (3) we have $Z \subs z_0$ by the test itself,
and $z_0 \subs z^\star$ by $[z_0] \subseteq [z^\star]$ and
\cref{lem:incl-implies-alu}; transitivity (\cref{lem:preorder}) gives
$Z \subs z^\star$.
\end{proof}

\begin{theorem}[Coverage monotonicity]\label{thm:coverage}
Let $A$ be the antichain at $\dk$ before an execution of \code{add} and $A'$
the antichain after.  Then $\Cov(A) \subseteq \Cov(A')$.  Consequently, for any
interleaving of workers, the covered region at each discrete state is
non-decreasing over time.
\end{theorem}

\FigCoverage

\begin{proof}
If \code{add} returns early at step (1) or (2), $A' = A$ and there is nothing to
show.  Otherwise write $R$ for the set of zones retired (in steps (3) and (4))
and $z^\star$ for the zone stored in step (5), so $A' = (A \setminus R) \cup
\{z^\star\}$.

By \cref{lem:replacement}, every $Z \in R$ satisfies $Z \subs z^\star$, i.e.\
$[Z] \subseteq \kappa([z^\star])$.  Applying $\kappa$ and using monotonicity and
idempotence,
\[
  \kappa([Z]) \;\subseteq\; \kappa\big(\kappa([z^\star])\big) \;=\;
  \kappa([z^\star]).
\]
Hence
\[
  \Cov(A) \;=\; \Cov(A\setminus R) \cup \bigcup_{Z \in R}\kappa([Z])
          \;\subseteq\; \Cov(A\setminus R) \cup \kappa([z^\star])
          \;=\; \Cov(A').
\]
Since stripe updates at a given $\dk$ are serialised by that stripe's exclusive
lock, the sequence of antichains at $\dk$ is totally ordered in time, and the
statement extends from one update to the whole execution by induction.
\end{proof}

\begin{corollary}[The shared-lock probe is sound]\label{cor:probe}
Suppose \code{probe\_kill} evaluates to true for a successor $z$ at $\dk$ at
time $t$, under a shared lock which is then released.  Then for every later time
$t' \ge t$ --- in particular, at the moment the worker decides to discard $z$
--- we still have $[z] \subseteq \Cov(A_{t'})$.  Discarding $z$ therefore never
discards a zone that the passed set does not cover.
\end{corollary}

\begin{proof}
\code{probe\_kill} returns true only after exhibiting a live $Z \in A_t$ with
$[z] \subseteq \kappa([Z])$, hence $[z] \subseteq \Cov(A_t)$.  Apply
\cref{thm:coverage} repeatedly to get $\Cov(A_t) \subseteq \Cov(A_{t'})$.
\end{proof}

\begin{remark}[what this does and does not license]
\cref{cor:probe} licenses reading the antichain without excluding writers.  It
does \emph{not} license reading it without any synchronisation: the probe still
takes the shared lock, because an exclusive writer may be halfway through
rewriting the chain's pointers and the reader must not observe a torn
structure.  The theorem is about the \emph{semantics} of the answer being
stable; the lock is about the \emph{representation} being intact.  Both are
needed, and conflating them is the classic way to get this wrong.
\end{remark}

\begin{remark}[why this is not free]
The property is a consequence of the specific update rule, not of antichains in
general.  An antichain algorithm that discarded a stored zone for \emph{any}
other reason --- a capacity bound, a heuristic eviction, a re-widening pass that
replaced a zone with a non-comparable one --- would break
\cref{lem:replacement} and with it the whole protocol.  This is worth stating in
the paper, because it identifies precisely what a competing design must
preserve.
\end{remark}

\subsection{The convex-union test is exact}

\begin{theorem}[Exactness of the convex-union merge]\label{thm:hull}
Let $A, B$ be canonical non-empty DBMs of dimension $\dim$ and let $H$ be their
entrywise maximum (the convex hull).  For $Z \in \{A,B\}$ call $(i,j)$, $i \ne
j$, a \emph{strict facet} of $Z$ if $Z[i][j] < H[i][j]$, and put
\[
  R_{Z,i,j} \;=\; [H] \,\cap\, \neg\big(x_i - x_j \prec_{Z} Z[i][j]\big).
\]
The procedure \code{dbm\_try\_union} accepts the merge iff
$R_{A,i,j} \subseteq [B]$ for every strict facet of $A$ and
$R_{B,i,j} \subseteq [A]$ for every strict facet of $B$.  Then:
\[
  \text{the procedure accepts}
  \quad\Longleftrightarrow\quad
  [H] \;=\; [A] \cup [B].
\]
In particular the merge never adds a valuation belonging to neither operand.
\end{theorem}

\begin{proof}
$[A]\cup[B] \subseteq [H]$ always, since $H$ is entrywise looser than both.

\emph{($\Rightarrow$) Accepting implies exactness.}
Let $u \in [H]$ and suppose $u \notin [A]$.  Because $A$ is canonical, $[A]$ is
exactly the set of valuations satisfying all constraints
$x_i - x_j \prec_A A[i][j]$, so there is a pair $(i,j)$ whose constraint $u$
violates; $i \ne j$, since a non-empty canonical DBM has $A[i][i] = (\le 0)$,
which every valuation satisfies.  If $A[i][j] = H[i][j]$ then $u \in [H]$ would
satisfy that very constraint, a contradiction; and $A[i][j] \le H[i][j]$ by
construction.  Hence $A[i][j] < H[i][j]$: it is a \emph{strict} facet.  Since
$u \in [H]$ and $u$ violates the constraint, $u \in R_{A,i,j}$, and acceptance
gives $R_{A,i,j} \subseteq [B]$, so $u \in [B]$.  Symmetrically for
$u \notin [B]$.  Therefore $[H] \subseteq [A]\cup[B]$.

\emph{($\Leftarrow$) Exactness implies accepting.}
Assume $[H] = [A]\cup[B]$ and let $(i,j)$ be a strict facet of $A$.  Every
$u \in R_{A,i,j}$ lies in $[H] = [A]\cup[B]$ and violates a constraint of $A$,
so $u \notin [A]$, hence $u \in [B]$.  Thus $R_{A,i,j} \subseteq [B]$, and the
same holds for $B$; the procedure accepts.
\end{proof}

\begin{remark}[faithfulness to the implementation]
The implementation constructs $R_{Z,i,j}$ by copying $H$, tightening the
\emph{reverse} entry $[j][i]$ with the negated bound --- $\neg(x_i - x_j < v)$
is $x_j - x_i \le -v$, and $\neg(x_i - x_j \le v)$ is $x_j - x_i < -v$ --- and
re-closing.  When that entry is not tighter than $H[j][i]$, the negation is
already implied by $H$ and $R = [H]$, which is the correct region.  When
$R = \emptyset$ the containment holds vacuously and the check is skipped.  So
the implemented test is exactly the one in \cref{thm:hull}.  A debug mode
(\code{PTCTLMC\_CHECKMERGE}) independently verifies every accepted merge with
exact federation subtraction --- $[H] \setminus ([A]\cup[B]) = \emptyset$ ---
which is a decision procedure, not a sanity check.
\end{remark}

\subsection{Deciding inclusion on the compressed representation}

\begin{lemma}[Minimal-graph inclusion]\label{lem:mingraph}
Let $G$ be a finite set of constraints $(i,j,b)$ whose closure is the canonical
DBM $Z'$, and let $z$ be canonical.  Then
\[
  [z] \subseteq [Z'] \quad\iff\quad z[i][j] \le b \ \text{ for every } (i,j,b) \in G .
\]
\end{lemma}

\begin{proof}
($\Rightarrow$) For canonical DBMs, $[z] \subseteq [Z']$ is equivalent to
$z[i][j] \le Z'[i][j]$ entrywise.  Since $Z'$ is the closure of $G$, we have
$Z'[i][j] \le b$ for each $(i,j,b) \in G$, so $z[i][j] \le b$.

($\Leftarrow$) Fix $i,j$.  $Z'[i][j]$ is the weight of a shortest path
$i = k_0 \to k_1 \to \cdots \to k_m = j$ using edges of $G$, say with bounds
$b_1,\dots,b_m$, so $Z'[i][j] = b_1 + \cdots + b_m$ in bound arithmetic.  Since
$z$ is canonical it satisfies the triangle inequality along that same path:
$z[i][j] \le z[k_0][k_1] + \cdots + z[k_{m-1}][k_m]$.  By hypothesis each
$z[k_{t-1}][k_t] \le b_t$, and bound addition is monotone in each argument, so
$z[i][j] \le b_1 + \cdots + b_m = Z'[i][j]$.  As $i,j$ were arbitrary,
$z \le Z'$ entrywise, i.e.\ $[z] \subseteq [Z']$.
\end{proof}

\begin{remark}
\cref{lem:mingraph} is the correctness statement of UDBM's
\code{dbm\_isSubsetEqMinDBM}, and we claim no novelty for the test itself; we
state and prove it because our use of it is what is new.  The test costs
$O(|G|)$, with $|G| \approx \dim$ in practice, against $O(\dim^3)$ for the
decompression a dense comparison would need.  It is what allows the entire
common case of insertion --- including the shared-lock probe of
\cref{fig:protocol} --- to run without ever materialising a stored zone.
\end{remark}

\subsection{Correctness of the parallel search}

\begin{theorem}[Soundness, completeness, termination]\label{thm:correct}
For every interleaving of the workers, \code{reach\_engine} terminates and
returns \emph{true} iff a state satisfying the target predicate is reachable.
\end{theorem}

\begin{proof}
\emph{Soundness.}  Every zone tested against the target is produced by the exact
zone algebra from a stored zone, starting at the exact initial zone: guard
conjunction, resets, delay and invariant conjunction are exact operations.  The
abstraction $\subs$ is used only to \emph{discard} candidates, never to enlarge
a zone, so no valuation is invented.  When widening is active (either forced,
or switched on adaptively) it can enlarge zones; \cref{prop:extra-sound} makes
that harmless provided the query's own clock constants are within $L,U$, which
the engine enforces by flooring $L$ and $U$ with them before the search starts
(\code{extend\_bounds}, and the per-state floors $L^q, U^q$).  Hence a returned
\emph{true} exhibits a genuinely reachable target state.

\emph{Completeness.}  We show by induction on $k$ the invariant

\begin{quote}
$(\mathrm{I}_k)$: at the barrier ending level $k$, for every symbolic state
$(\dk, Z)$ reachable in at most $k$ symbolic steps, there is a \emph{live}
stored node $(\dk, Z^\star)$ with $Z \subs Z^\star$; and every live node has
either been expanded already or belongs to the frontier of level $k$.
\end{quote}

$(\mathrm{I}_0)$ holds: the initial state is inserted and is the frontier.

Assume $(\mathrm{I}_k)$ and let $(\dk', W)$ be reachable in $k{+}1$ steps, via a
predecessor $(\dk, Z)$ reachable in $k$ steps.  By $(\mathrm{I}_k)$ there is a
live $Z^\star$ with $Z \subs Z^\star$, expanded at some level $\le k{+}1$.
Because $\subs$ is a simulation for the symbolic successor operator
(\cref{prop:alu-finite}), expanding $Z^\star$ produces some $(\dk', W^\star)$
with $W \subs W^\star$, and that successor is offered to \code{add}.  Three
outcomes:
\begin{enumerate}[leftmargin=1.6em, itemsep=1pt]
\item it is dropped by the probe or by step (1)/(2).  Then some live $Y$ has
      $W^\star \subs Y$ at the moment of the test; by \cref{cor:probe} and
      \cref{thm:coverage} the covered region only grows, so at the barrier there
      is still a live $Y'$ with $W^\star \subs Y'$, and transitivity gives
      $W \subs Y'$.
\item it is stored and survives to the barrier.  Take $Y' = W^\star$.
\item it is stored and later retired within the same level.  By
      \cref{lem:replacement} the retiring insertion stored, at the same $\dk'$
      and at the same level, a zone subsuming it; iterating (the level performs
      finitely many insertions) yields a live $Y'$ at the barrier with
      $W^\star \subs Y'$.
\end{enumerate}
In all three cases $W \subs Y'$ for a live $Y'$, which is the first conjunct of
$(\mathrm{I}_{k+1})$.

For the second conjunct: a node stamped level $k{+}1$ that is live at the
barrier is visited by \code{for\_level} for its $\dk'$, and $\dk'$ is in some
worker's \emph{touched} set because an insertion at $\dk'$ succeeded during the
level; by \cref{lem:bucket} the bucketed union sees every such key.  So it
enters the frontier of level $k{+}1$ and will be expanded.  A node stamped an
earlier level was, by $(\mathrm{I}_k)$, already expanded.  Tombstones are never
revived, so no live node can have skipped its collection.

Finally, if a reachable target state exists, it is reachable in some finite
number of steps $k$, and $(\mathrm{I}_{k-1})$ guarantees a stored zone
subsuming its predecessor is expanded; the successor generated from it
satisfies the target because the target predicate is monotone with respect to
$\subs$ on the fragment the engine accepts (location and integer conditions are
determined by $\dk$; clock conditions are floored into $L,U$).  Hence the
target test fires and \code{found} is set.

\emph{Termination.}  In subsumption mode, \cref{prop:alu-finite} bounds the
number of pairwise non-subsumed zones per discrete state, and there are finitely
many discrete states (the packed key is a bijection onto a finite product of
finite ranges); so the passed set is finite and BFS stops.  In widening mode the
image of $\extraLU$ is finite (\cref{prop:extra-sound}), giving the same
conclusion.  The mixed case is \cref{thm:adapt}.
\end{proof}

\begin{lemma}[Bucketed union]\label{lem:bucket}
Let $S_1,\dots,S_p$ be the per-worker touched sets and $h$ a function into
$\{0,\dots,\NB_{\mathrm{buk}}-1\}$.  Deduplicating each bucket
$\{\, x \in \bigcup_t S_t : h(x) = b \,\}$ independently and concatenating
yields exactly the set $\bigcup_t S_t$, up to order.
\end{lemma}

\begin{proof}
$h$ is a function, so each $x$ lands in exactly one bucket, and that bucket
receives every copy of $x$ contributed by any worker.  Deduplication within the
bucket therefore emits $x$ exactly once, and no bucket emits an $x$ it did not
receive.  The concatenation is a partition of the union into $h$-classes.
\end{proof}

\begin{remark}
Order is not part of the specification: the previous serial union already
produced a schedule-dependent order (a dynamic-schedule loop appending into
per-thread buffers), and the search's result does not depend on frontier order.
What the merge \emph{is} sensitive to is order, which is why the verdict, not
the stored set, is the invariant --- see \cref{thm:determinism}.
\end{remark}

\begin{theorem}[Verdict determinism, state-count nondeterminism]
\label{thm:determinism}
The verdict is independent of the thread schedule.  The set of stored symbolic
states, and hence the reported state count, is not.
\end{theorem}

\begin{proof}
The first claim is immediate from \cref{thm:correct}: for every interleaving the
returned value equals ``a target state is reachable'', a property of the model
and query alone.

For the second, it suffices to observe that step (4) of \code{add} is
order-sensitive: whether $Z_1$ and $Z_2$ can be absorbed into the running hull
depends on the order in which they are offered, and that order depends on which
worker reached the stripe first.  Two schedules can therefore terminate with
different --- both valid --- antichains.  At one thread the engine takes the
\code{single} path, which skips the locks and the bucketed union entirely and
reproduces the sequential insertion order exactly, so single-threaded runs are
reproducible bit-for-bit.
\end{proof}

\begin{remark}[why this must be said out loud]
Reviewers and users will run the tool twice at $32$ threads and see two
different ``states explored'' numbers.  That is expected and benign, and the
tool should say so.  What must \emph{not} vary is the verdict, and
\cref{thm:correct} is exactly the statement that it does not.
\end{remark}

\subsection{The adaptive abstraction switch}

The engine starts in subsumption mode (exact zones, $\subs$ pruning) and
monitors the ratio of live zones to distinct discrete states at each level
boundary.  A ratio near $1$ means subsumption plus merging is collapsing the
antichains and widening would only cost precision; a ratio that climbs means the
antichains are filling with mutually incomparable zones, which pruning cannot
remove and widening can.  Above a threshold (default $4$) the engine switches to
$\extraLU$ for all \emph{subsequently generated} zones and never switches back.
The passed set is then mixed: exact zones stored before the switch, widened
zones after.

\begin{theorem}[Soundness and termination of the mixed mode]\label{thm:adapt}
Switching on $\extraLU$ at a level boundary, one-way, preserves
\cref{thm:correct}.
\end{theorem}

\begin{proof}
\emph{Soundness.}  After the switch, target tests are applied to widened zones.
Since $L$ and $U$ are floored by the query's clock constants before the search
begins, \cref{prop:extra-sound} gives that $\extraLU(Z)$ satisfies the target's
clock condition iff $Z$ does; the location and integer parts of the target are
functions of $\dk$, which widening does not touch.  So no false positive is
introduced.

\emph{Completeness.}  Two mixed situations arise.  (a) A stored \emph{exact}
zone fails to subsume an incoming \emph{widened} zone.  Then the widened zone is
stored as well: the only consequence is an extra stored state, and the
invariant $(\mathrm{I}_k)$ of \cref{thm:correct} is preserved a fortiori.  (b) A
stored \emph{widened} zone $W$ subsumes an incoming zone $Z$, so $Z$ is dropped.
This is sound because $\subs$ is a simulation as a relation between \emph{sets
of valuations} (\cref{prop:alu-finite}); the property does not refer to how $W$
was obtained.  Also $Z \subseteq \extraLU(Z)$, so widening a successor can only
help it be subsumed, never cause a needed state to be missed.

\emph{Termination.}  The switch occurs at a level boundary, hence after finitely
many levels.  Each level stores finitely many nodes (level $0$ is a singleton
and every symbolic state has finitely many successors), so the set $S_0$ of
zones stored before the switch is finite.  After the switch, every newly stored
zone lies in the image of $\extraLU$, which is finite
(\cref{prop:extra-sound}), and there are finitely many discrete states; so at
most finitely many further zones can be stored.  The total stored set is
therefore contained in $S_0 \cup (\text{finite})$, and BFS terminates.  The
one-way nature of the switch is essential: a mode that could switch back
arbitrarily often would not admit this argument.
\end{proof}

\begin{remark}
The adaptive switch is not a tuning heuristic bolted onto a fixed algorithm; it
is a decision between the two \emph{different} finiteness mechanisms of
\cref{sec:bg-finite}, taken from a signal the engine already computes for free.
Measured effect: on \code{single}, widening takes the stored count from $3873$
to $3676$ (\uppaal{} stores $3673$); on \code{phil} and \code{viking}, where the
ratio stays at $1.11$ and $1.00$, it changes no state count and would have cost
$20$--$35\%$ in time, so it correctly never fires.
\end{remark}

\subsection{Reclamation: amortisation and a soundness obligation}

\begin{proposition}[Amortised GC]\label{prop:gc}
Let the interner sweep whenever its allocated-record count reaches
$\mathit{target}$, and afterwards set $\mathit{target} := 2 \cdot
\mathit{live}$, where $\mathit{live}$ is the surviving count.  Then the total
sweep work over a run is $O(1)$ amortised per interned zone, and the allocated
count never exceeds twice the live count measured at the previous sweep.
\end{proposition}

\begin{proof}
A sweep costs $\Theta(a)$ where $a$ is the allocated count at the trigger, since
it walks the arena nodes and the hash slots once.  Between the sweep that set
the target to $2\ell$ and the next trigger at $a = 2\ell$, at least $\ell$ fresh
records were allocated (the count was $\ell$ immediately after the sweep).  So
the $\Theta(2\ell)$ cost is charged against $\ell$ allocations, i.e.\ $O(1)$
each.  The memory bound is the trigger condition itself.
\end{proof}

\begin{remark}
The predecessor of this trigger was $\mathit{allocated} > 50\,000$.  Because
the allocation counter is a high-water mark that never shrinks, that predicate
latched true at the first large level and then ran a full mark-and-sweep on
\emph{every} subsequent level.  Both halves of the sweep are serial and grow
with the state space while the expansion phase does not, so on
\code{arbiter17} it was $2.8$\,s of a $9.3$\,s run ($30\%$) and was what capped
the $12$-thread speedup at $\approx 5\times$.  The lesson generalises: in a
level-synchronous parallel search, any maintenance task whose cost grows with
the state space must be given an amortised trigger, or Amdahl's law will find
it.
\end{remark}

\begin{proposition}[Cache invalidation is a soundness obligation]
\label{prop:gc-cache}
Suppose the interner recycles freed record slots and freed indices in place
(rather than renumbering), and suppose each worker keeps a
$\mathit{hash} \mapsto \mathit{index}$ cache whose hits are validated by
comparing record bytes.  Then clearing every worker cache during the sweep is
\emph{necessary and sufficient} for the interner to remain sound.
\end{proposition}

\begin{proof}
\emph{Necessity.}  Sweeping does not erase a freed record; it pushes the memory
onto a free list.  A freed record therefore still compares byte-equal to the
zone it used to hold, so a stale cache line for that zone passes validation and
returns an index that is on the free list.  The caller stores that index in a
passed-set node.  The next interning that pops the same index writes a different
record into it, and the node now denotes a \emph{different} zone: the state
keeps its discrete key and acquires an unrelated zone, so the search explores
states the model cannot reach.  (This is not hypothetical; it was observed on
\code{fischer6} and larger, where it turned a satisfied mutual-exclusion
property into a spurious counterexample.)  Byte equality cannot distinguish the
two cases, so no cheaper validation suffices.

\emph{Sufficiency.}  After clearing, no cache line can return a pre-sweep index,
so every interning goes to its shard under the shard lock and consults the
rebuilt hash table, which by construction contains exactly the survivors.
Indices already held by passed-set nodes and by the frontier remain valid,
because the mark phase marks every index reachable from a live arena node or
from the next frontier, and only unmarked indices are freed.  The sweep runs
between levels with all workers joined, so there is no other holder of an index
at that moment.
\end{proof}

\begin{proposition}[Shard and slot indices must be independent]\label{prop:shard}
Let a hash table of $m$ slots be split into $2^b$ shards, let the shard be
selected by the low $b$ bits of the hash, and let each shard's slot index also
be taken from the low bits.  Then every key routed to shard $s$ has the same
residue modulo $2^b$, so at most $m/2^b$ of the shard's slots can be a home
position, and the expected probe length at load factor $\alpha$ degrades from
$\Theta\big(1/(1-\alpha)\big)$ to $\Theta(2^b \alpha)$.
\end{proposition}

\begin{proof}
Write $h$ for the hash.  The shard is $h \bmod 2^b$ and the home slot is
$h \bmod m$ with $2^b \mid m$.  Then
$(h \bmod m) \bmod 2^b = h \bmod 2^b = s$ is constant across the shard, so home
slots are confined to the arithmetic progression $\{\,s, s+2^b, s+2\cdot 2^b,
\dots\,\}$, of size $m/2^b$.  With $\alpha m$ entries in the shard, the expected
number occupying a given home slot's probe run is $\Theta(\alpha m / (m/2^b)) =
\Theta(2^b\alpha)$.
\end{proof}

\begin{remark}
This is elementary once stated, and it cost a full sharding rewrite to
discover: on \code{arbiter16} the sharded table did $31.2$ probes per locked
lookup where the unsharded one did $1.68$, and each probe chases a chunk pointer
and byte-compares a record.  The fix is to mix the hash and take the \emph{high}
bits for the shard, making the two indices independent.
\end{remark}
